\documentclass{article}
\usepackage[T1]{fontenc}
\usepackage{lmodern}
\usepackage{graphicx}
\usepackage{amsmath,amssymb}
\usepackage[a4paper,margin=2cm]{geometry}
\usepackage[absolute,overlay]{textpos}
\setlength{\TPHorizModule}{1mm}
\setlength{\TPVertModule}{1mm}
\usepackage[indonesian]{babel}
\usepackage{hyperref}
\setlength{\emergencystretch}{2em}
\title{02-Landasan-Matematika-Kriptografi-(2026)}
\author{pdf\_ppt\_to\_latex}
\date{}

\begin{document}

% --- unit llm 1 ---
% === Halaman 1 (llm) ===

\begin{center}
{\large IF4020 Kriptografi}

\vspace{20mm}
{\huge \textcolor{red}{02 -- Landasan Matematika untuk Kriptografi}}

\vspace{15mm}
Oleh: Rinaldi M

\vspace{10mm}
\textbf{Program Studi Teknik Informatika}\
\textbf{Sekolah Teknik Elektro dan Informatika}\
\textbf{Institut Teknologi Bandung}\
\textbf{2026}
\end{center}

\newpage

% --- unit llm 2 ---
% === Halaman 2 (llm) ===

\section*{Materi Matematika untuk Kriptorafi:}

\begin{enumerate}
    \item Teori Bilangan
    \begin{itemize}
        \item Integer dan aritmetika modulo
        \item Algoritma euclidean
        \item Kekongruenan
        \item Relatif prima
        \item Balikan (invers) modulo
        \item Bilangan prima
    \end{itemize}
    \item Probabilitas dan Statistik
    \item Kompleksitas algoritma
    \item Teori Informasi
    \item Aljabar abstrak
\end{enumerate}

\vspace{10mm}

\begin{tabular}{|p{9cm}|}
\hline
No. 1 s/d 3 sudah dipelajari di dalam kuliah Matematika Diskrit dan Probabilitas dan Statistik \\
\hline
No 5 akan dibahas pada materi $ECC$ (Elliptic Curve Cryptography) \\
\hline
\end{tabular}

\newpage

% --- unit llm 3 ---
% === Halaman 3 (llm) ===

Contoh:
\begin{enumerate}[(i)]
  \item $23 \pmod{5} = 3$ karena $23$ dibagi $5 = 4 + \text{ sisa } 3$
  \item $-41 \pmod{9} = 4$ karena $|-41| \pmod{9} = 5$, dan $9 - 5 = 4$
  \item $17 \equiv 2 \pmod{3}$ karena $3 \mid (17 - 2)$
  \item $-7 \equiv 15 \pmod{11}$ karena $11 \mid (-7 - 15)$
  \item $23$ dan $40$ relatif prima sebab $\text{PBB}(23, 4) = 1$
  \item $4^{-1} \pmod{9} \equiv 7 \pmod{9}$ karena $4 \cdot 7 \equiv 1 \pmod{9}$
  \item $23^{-1} \pmod{10} = -3 \pmod{10}$ karena $23 \cdot (-3) \equiv 1 \pmod{10}$
\end{enumerate}

Latihan: (a) Hitung $-24 \pmod{11} = ?$
(b) $12^{-1} \pmod{5} \equiv ?$

\newpage

% --- unit llm 4 ---
% === Halaman 4 (llm) ===

\section*{Teori Informasi (Information Theory)}

\begin{itemize}
    \item \textit{Information theory}: cabang ilmu yang mempelajari \textbf{kuantisasi, penyimpanan, transmisi, dan pengolahan informasi}
    \item Contoh kuantisasi:
    \begin{itemize}
        \item 1 bit untuk mengkodekan jenis kelamin (M/F)
        \item 3 bit untuk mengkodekan nama hari (ada 7 hari)
        \item 4 bit untuk mengkodekan 0 s/d 9
    \end{itemize}
    \item Contoh penyimpanan: pemampatan data untuk mengurangi ukuran ruang \textit{storage}
\end{itemize}

\newpage

% --- unit llm 5 ---
% === Halaman 5 (llm) ===

\begin{itemize}
    \item Salah satu metrik penting di dalam teori informasi adalah \textbf{entropi}.
    \item \textbf{Entropi} mengukur \textit{ketidakpastian} atau jumlah rata-rata informasi di dalam pesan.
    \item Semakin acak suatu data, semakin tinggi entropinya. Cipherteks adalah pesan dengan entropi yang tinggi.
    \item Entropi biasanya dinyatakan dalam satuan bit.
    \item Secara umum, entropi pesan $X$ dihitung dengan rumus:
\end{itemize}

\vspace{5mm}

\[ H(X) = - \sum_{i=1}^{n} p(x_i) \log_2 p(x_i) \]

\vspace{5mm}

$X$ = variabel acak yang menyatakan pesan \\
$x_i$ = simbol ke-$i$ di dalam pesan \\
$n$ = banyak simbol di dalam pesan \\
$p(x_i)$ = peluang kemunculan $x_i$

\newpage

% --- unit llm 6 ---
% === Halaman 6 (llm) ===

\begin{itemize}
  \item Contoh: misalkan pesan $X = \text{'AABBCBDB'}$
  \begin{align*}
    n &= 4 \text{ (yaitu huruf A, B, C, D)}\
    p(A) &= 2/8, p(B) = 4/8\
    p(C) &= 1/8, p(D) = 1/8
  \end{align*}
\end{itemize}

\begin{align*}
  H(x) &= -\{2/8 \, ^2\log(2/8) + 4/8 \, ^2\log(4/8) + 1/8 \, ^2\log(1/8) + 1/8 \, ^2\log(1/8) \}\
  &= -\{1/4 \, ^2\log(1/4) + 1/2 \, ^2\log(1/2) + 1/8 \, ^2\log(1/8) + 1/8 \, ^2\log(1/8) \}\
  &= -\{(1/4) \, (-^2\log(4) + (1/2) \, (-^2\log(2) + (1/8) \, (-^2\log(8) + (/8) \, (-^2\log(8) \}\
  &= -\{(1/4) \, (- 2) + (1/2) \, (-1) + (1/8) \, (-3) + (1/8)(-3) \}\
  &= - \{-1/2 - 1/2 - 3/8 - 3/8)\
  &= - (-1.75)\
  &= 1.75
\end{align*}

Entropi $= 1,75$ bit per simbol

\newpage

% --- unit llm 7 ---
% === Halaman 7 (llm) ===

\begin{itemize}
    \item Nilai entropi berkisar dalam rentang:
    \[ 0 \le H(X) \le \log_2(n) \]
    \item Entropi minimum $= 0$ bit
    \begin{itemize}
        \item $\to$ Tidak ada ketidakpastian
        \item $\to$ Terjadi jika hanya ada satu simbol dengan peluang $1$ (pasti).
    \end{itemize}
    Contoh: $n = 1$, misalkan A dan $p(A) = 1 \to H(X) = -1 \cdot \log_2 1 = 0
    \item Entropi maksimum $= \log_2(n)$ bit
    \begin{itemize}
        \item $\to$ Terjadi jika semua simbol muncul dengan peluang sama, yaitu $1/n$.
        \item $\to$ Ketidakpastian paling tinggi, setiap simbol membawa informasi maksimal.
    \end{itemize}
    Contoh: Empat simbol dengan peluang sama ($0.25$) $\to H(X) = -\sum_{i=1}^{4} 0.25 \log_2 0.25 = 2$ bit
\end{itemize}

\newpage

% --- unit llm 8 ---
% === Halaman 8 (llm) ===

\begin{itemize}
    \item Jika plainteks/cipherteks hanya menggunakan 26 huruf alfabet (A--Z), maka entropi Shannon maksimum terjadi ketika seluruh 26 huruf memiliki probabilitas yang sama:
    \[ p_i = \frac{1}{26} \]
    sehingga
    \begin{align*}
    H_{max} &= - \sum_{i=1}^{26} p_i \log_2 p_i \\
    &= - \sum_{i=1}^{26} \frac{1}{26} \log_2 \frac{1}{26} = \log_2 26 = 4.7004 \text{ bit/karakter}
    \end{align*}
    Jadi, untuk teks yang hanya terdiri dari 26 huruf alfabet, entropi Shannon maksimum adalah sekitar 4,7004 bit/karakter.
    
    \item Jika pesan terdiri dari 256 karakter ASCII, maka entropi maksimum adalah $\log_2 256 = 8$
    
    \item Makin besar entropi, makin sulit memecahkan cipherteks.
\end{itemize}

\newpage

% --- unit llm 9 ---
% === Halaman 9 (llm) ===

Plainteks: (P)
\begin{center}
\fbox{
\begin{minipage}{0.6\textwidth}
Dinas Pendidikan Kota Ternate meminta kepada pihak sekolah dan orang tua siswa untuk jenjang pendidikan SD dan SMP se-Kota Ternate untuk melarang para siswa membawa permainan lato-lato yang sedang tren itu ke sekolah, karena akan mengganggu kegiatan belajar mengajar yang dinilai berbahaya sehingga mengantisipasi kecelakaan bagi anak di daerah itu.
\end{minipage}
}
\end{center}
\vspace{5mm}
$H(P)=3,8988$ bit/karakter
\vspace{10mm}
Cipherteks: (C)
\begin{center}
\fbox{
\begin{minipage}{0.6\textwidth}
HAWFHZDOHAHANGOMKLGFCVWFLBOCPRKFGHNOFNIN\\
SPGNLHMPFVBEFWMVFWTBWRHSFZRWK FQMVHPAF WIK\\
DOHAHANGPFEWNF PFNKLHMPFGLBAMGF C XK FQ MOCFV\\
AVKANIAVHSFZR NCODG AF OHYUVG W G F O F G X E G F M L F W I T\\
HEFWTBVCPAYQOBLHMPFVBPVADOVGNWFG NOV C A R K V A\\
RQOBRGGF WH C V F G V Y Y U D O R V G V Y C F W A B A P V S V G C H G C I\\
DRFIFH B A P V Q R V O C K A F W S G S H S N I H S O B D H F V N G F W D G R G R F\\
W G N H A N F C V I D G S W Z
\end{minipage}
}
\end{center}
\vspace{5mm}
$H(C)=4,3035$ bit/karakter

\vspace{71.28mm}

Entropi cipherteks meningkat

Jadi cipherteks ini mencapai entropi sekitar:
\[ \frac{4,3035}{4,7004} \times 100\% \approx 91,55\% \]
dari entropi maksimum pesan 26 huruf alfabet

% deskripsi: Tanda panah biru yang mengarah ke bawah, menghubungkan nilai H(P) dengan H(C).
% bbox normalisasi: [0.8300, 0.2000, 0.8600, 0.4400]
\begin{textblock*}{6.30mm}(174.30mm,59.40mm)
\noindent\includegraphics[width=6.30mm,height=71.28mm]{images/llm/page_009_img_01.png}%
\end{textblock*}

\end{document}
